\documentclass[10pt,a4paper]{amsart}
\usepackage[margin=19mm]{geometry}
\usepackage{amsmath,amssymb,mathtools}
\usepackage[scr=rsfs,cal=euler]{mathalfa}
\usepackage{xfrac}
\usepackage[unicode,hidelinks,bookmarks=false]{hyperref}
\newcommand{\ie}{i.e.\ }
\newcommand{\cf}{cf.\ }
\newcommand{\resp}{respectively\ }
\newcommand{\Subsection}{\subsection}
\providecommand{\nobreakdash}{\nobreak\hbox{-}}
\setlength{\emergencystretch}{2em}
\multlinegap0pt
\usepackage{xcolor}
\usepackage{amssymb}
\usepackage{enumerate}
\usepackage{bm}
\usepackage{float}
\usepackage{tikz}
\newtheorem{theorem}{Theorem}
\title{A Mathematical Model of Dengue Transmission Incorporating Hospital Capacity and Threshold-Based Fogging Interventions}
\author{Dipo Aldila, Joseph P\'aez Ch\'avez, Aytül Gökçe, Thomas Götz, Burcu Gürbüz$^*$}
\date{Source: arXiv:2607.18140v1 (CC BY 4.0)}
\begin{document}
\maketitle

\noindent\emph{Source and licence.} A Mathematical Model of Dengue Transmission Incorporating Hospital Capacity and Threshold-Based Fogging Interventions. Version arXiv:2607.18140v1 (CC BY 4.0), distributed by arXiv under the Creative Commons Attribution 4.0 International licence, http://creativecommons.org/licenses/by/4.0/, as recorded in the arXiv licence metadata for this exact version and shown on the abstract page https://arxiv.org/abs/2607.18140v1. Full licence notice: \texttt{licenses/arxiv-2607.18140-CC-BY-4.0.txt}.\par\smallskip

\noindent\textit{Editorial scope.} Counted unit: the article's theorem theorem (source0.tex lines 481--483). The article's complete original proof of it is delivered with it (source0.tex lines 485--517, one \texttt{proof} environment of 33 lines). No mathematics is written, repaired or re-derived: the statement and the proof are the article's own lines, in the article's own notation, and no retained line is substituted or re-flowed.  Retained uncounted as the source-local context the proof needs: lines 372--381; lines 439--449.  Nothing was omitted from any retained span, so the package carries no omission record.  No retained line contains a citation, so the wrapper carries no bibliography and no bibliography entry.  No retained line contains a figure environment, an image-inclusion command or drawing code, so no image is delivered and the package declares no asset dependency.  Editorial formatting only, in addition: an AMS \texttt{amsart} preamble that restates the article's own declarations and macros used by the retained lines, namely the environments \texttt{theorem} on separate counters, together with the standard packages those lines need.  The retained environments print in this document as Theorem 1 in order of appearance, whereas the article prints the same items with the numbering of its own full document.  The complete native source of the article as distributed in the arXiv e-print for this exact version is bundled unchanged as \texttt{sources/205632/source0.tex}.
\begin{equation}\label{model}
    \begin{aligned}
        \dot{S} &= \Lambda_h - \beta_h S V - \mu_h S,  \\
    \dot{A} &= p \beta_h S V - \left(\alpha + \gamma_0 + \mu_h\right)A,  \\
    \dot{I} &= (1-p)\beta_h S V + \alpha A - f(I) - (\mu_h + \delta)I,  \\
    \dot{R} &= \gamma_0 A + f(I) - \mu_h R, \\
    \dot{U} &= \Lambda_m - g(U,A,I) - (\mu_m + h(I))U,  \\
    \dot{V} &= g(U,A,I) - (\mu_m + h(I))V,
    \end{aligned}
\end{equation}

\subsection{Case when $I(t)< kC$}
When the number of infected individuals are less than the first threshold ($kC$, with $k \in (0,1)$), model~\eqref{model} read as follows:
\begin{eqnarray}\label{modela}
    \dot{S} &=& \Lambda_h - \beta_h S V - \mu_h S, \nonumber \\
    \dot{A} &=& p \beta_h S V - \left(\alpha + \gamma_0 + \mu_h\right)A, \nonumber \\
    \dot{I} &=& (1-p)\beta_h S V + \alpha A - \gamma_1 I - (\mu_h + \delta)I, \nonumber \\
    \dot{R} &=& \gamma_0 A + \gamma_1 I - \mu_h R, \\
    \dot{U} &=& \Lambda_m - \beta_{m1}UA - \mu_m U, \nonumber \\
    \dot{V} &=& \beta_{m1}UA - \mu_m V. \nonumber
\end{eqnarray}
For any non-negative initial condition of $S, A, I, R, U,$ and $V$, it can be shown that the direction of vector field in the boundary region of $\mathbb{R}_+^6$ always has an inward direction. This means that all solutions will always remain positive for any $t>0$.

\begin{theorem}
    Dengue-free equilibrium is locally asymptotically stable when $\mathcal{R}_0<1$ and unstable if $\mathcal{R}_0>1$.
\end{theorem}

\begin{proof}
After substituting the Dengue-free equilibrium $\mathcal{E}_0$ into the Jacobian matrix for the system~\eqref{modela}, we calculate the following six eigenvalues
\begin{align*}
    \lambda_{1,2}(\mathcal{E}_0) &= -\mu_h\;,\\
    \lambda_3(\mathcal{E}_0) &= -\mu_m\;,\\
    \lambda_4(\mathcal{E}_0) &= -(\gamma_1+\mu_h+\delta)\;, \\
    \lambda_{5,6}(\mathcal{E}_0) &= -\frac{1}{2} \left( -(\alpha+\gamma_0+\mu_h+\mu_m) \pm \sqrt{(\alpha+\gamma_0+\mu_h-\mu_m)^2+4p\beta_h \beta_{m1}\frac{\Lambda_h\Lambda_m}{\mu_h\mu_m}} \right)\;.
\end{align*}
It is easy to see, that all eigenvalues are negative provided $p\beta_h \beta_{m1}\frac{\Lambda_h\Lambda_m}{\mu_h\mu_m} < (\alpha+\gamma_0+\mu_h)\mu_m$, i.e.~if $\mathcal{R}_0<1$.

%
%    After substituting the Dengue-free equilibrium, $\mathcal{E}_0$, into the Jacobian matrix for the system \ref{modela}, we calculate the eigenvalues of the resulting matrix. The eigenvalues are
%    \begin{align}
%     &-\mu_{h}, -\mu_{h}, -\mu_{m}, -(\gamma_{1}+\mu_{h}+\delta), \nonumber \\
%    &\frac{1}{{2 \mu_{h} \mu_{m}}} (-\mu_{h} \mu_{m} \left(\alpha+\gamma_{0}+\mu_{h}+\mu_{m}\right)+\sqrt{d}), \nonumber \\
%    &\frac{-1}{{2 \mu_{h} \mu_{m}}} (\mu_{h} \mu_{m} \left(\alpha+\gamma_{0}+\mu_{h}+\mu_{m}\right)+\sqrt{d}), \nonumber \\
%    \text{where} \nonumber \\
%    &d = \mu_{m} \mu_{h} \left(\mu_{h}^{3} \mu_{m}+2 \mu_{m} \left(\alpha -\mu_{m}+\gamma_{0}\right) \mu_{h}^{2}+\mu_{m} \left(\alpha -\mu_{m}+\gamma_{0}\right)^{2} \mu_{h}+4 p \Lambda_{h} \Lambda_{m} \beta_{h} \beta_{\mathit{m1}}\right). \nonumber
%    \end{align}
%
%
%    \textcolor{blue}{A.G.\\ .\hspace{1.5cm}$d = \mu_{m} \mu_{h} \left(\mu_{h} \mu_{m}^{3}-2 \mu_{m}^{2} \mu_{h}\left(\alpha +\mu_{h}+\gamma_{0}\right) +\mu_{m} \mu_{h} \left(\alpha +\mu_{h}+\gamma_{0}\right)^{2}+4 p \Lambda_{h} \Lambda_{m} \beta_{h} \beta_{\mathit{m1}}\right)$\\\\
%    I believe all three results are same, just interpreted in different ways. }

%    \textcolor{orange}{TG: I get for the last two eigenvalues
%    \begin{equation*}
%        \frac{1}{2} \left( -(\alpha+\gamma_0+\mu_h+\mu_m) \pm \sqrt{(\alpha+\gamma_0+\mu_h-\mu_m)^2+4p\beta_h %\beta_{m1}\frac{\Lambda_h\Lambda_m}{\mu_h\mu_m}} \right)
%    \end{equation*}
%    In this case, the signs are not directly obvious -- at least to me at first sight... Can someone please check a third time and ome up with something even different??}
%
%
%    It is evident that all the eigenvalues are negative, except for $\lambda_5$. After some calculations, we determine that it is negative for $\mathcal{R}_0<1$. Since all the eigenvalues are negative, this concludes the proof.
\end{proof}

\end{document}
